\documentclass[10pt, a4paper]{article}

% --- Encoding and language ---
\usepackage[T1]{fontenc}
\usepackage[utf8]{inputenc}
\usepackage[english]{babel}
\usepackage[warn]{textcomp}

% --- Page geometry (two-column) ---
\usepackage{geometry}
\geometry{a4paper, left=15mm, right=15mm, top=20mm, bottom=20mm}

% --- Math and fonts ---
\usepackage{amsmath}
\usepackage{amssymb}
\usepackage{amsthm}

% --- Figures and tables ---
\usepackage{graphicx}
\usepackage{xcolor}
\usepackage{booktabs}
\usepackage{tabularx}
\usepackage{enumitem}
\usepackage{float}

% --- START OF FOOTER SECTION ---
\usepackage{fancyhdr}
\usepackage[colorlinks=true, urlcolor=blue, linkcolor=blue]{hyperref}
\usepackage[numbers]{natbib}

\pagestyle{fancy}
\fancyhf{}
\renewcommand{\headrulewidth}{0pt}
\renewcommand{\footrulewidth}{0pt}
\lfoot{\scriptsize \href{https://willrg.com}{https://willrg.com}}
\cfoot{\scriptsize \thepage}
\rfoot{\scriptsize \href{https://doi.org/10.5281/zenodo.17115270}{DOI: 10.5281/zenodo.17115270}}

\fancypagestyle{plain}{
  \fancyhf{}
  \renewcommand{\headrulewidth}{0pt}
  \renewcommand{\footrulewidth}{0pt}
  \lfoot{\scriptsize \href{https://willrg.com}{https://willrg.com}}
  \cfoot{\scriptsize \thepage}
  \rfoot{\scriptsize \href{https://doi.org/10.5281/zenodo.17115270}{DOI: 10.5281/zenodo.17115270}}
}
% --- END OF FOOTER SECTION ---

% --- TikZ ---
\usepackage{tikz}
\usetikzlibrary{shapes, positioning, arrows.meta, calc}

% --- tcolorbox for styled theorems ---
\usepackage[most]{tcolorbox}
\tcbuselibrary{theorems, breakable, skins}

% --- Theorem environments ---
\newtheorem{theorem}{Theorem}[section]
\newtheorem{lemma}[theorem]{Lemma}
\newtheorem{corollary}[theorem]{Corollary}
\newtheorem{definition}[theorem]{Definition}
\newtheorem{proposition}[theorem]{Proposition}
\newtheorem{remark}[theorem]{Remark}
\newtheorem{principle}[theorem]{Principle}

% --- For two-column layout ---
\usepackage{multicol}
\usepackage{cuted}
\usepackage{stfloats}

% =============================================================================
% DOCUMENT
% =============================================================================

\begin{document}

% --- Title and Abstract (full-width) ---
\begin{center}
{\Large\bfseries WILL Part~I: Relational Geometry}\\[8pt]
{\itshape A Foundational Framework for the Unification of\\
Special Relativity, General Relativity, and Orbital Mechanics\\
from a Single Ontological Principle}\\[12pt]
{Anton Rize\thanks{This work is archived on Zenodo: \href{https://doi.org/10.5281/zenodo.17115270}{DOI:10.5281/zenodo.17115270}}}\\
{\small antonrize@willrg.com}\\[4pt]
{\small October 2025}
\end{center}

\vspace{6pt}

\begin{center}
\begin{minipage}{0.95\textwidth}
\small
\textbf{Abstract.}
This paper, the first in the \textbf{WILL Trilogy}
\hypertarget{trilogy}{}\label{trilogy}
(\href{https://willrg.com/documents/WILL_RG_I.pdf}{Part~I: Relational Geometry};
\href{https://willrg.com/documents/WILL_RG_II.pdf}{Part~II: Relational Approach to Dark Sector Problem in Cosmology};
\href{https://willrg.com/documents/WILL_RG_III.pdf}{Part~III: Relational Quantum Geometry}),
applies extreme methodological constraints and establishes \textit{Relational Geometry (RG)}: a foundational framework where spacetime is an emergent property of relational energy transformations. This shift establishes an ontological transition from \textit{descriptive} to \textit{generative} physics: instead of introducing laws to model observations, it derives them as necessary consequences of RG itself---turning physics from a catalogue of phenomena into the logical unfolding of inevitable geometrical constraints on closed relational carriers $S^1$ (directional) and $S^2$ (omnidirectional).

Without metrics, tensors, or free parameters, it reproduces Lorentz factors, the energy--momentum relation, Schwarzschild and Einstein field equations via the dimensionless projections $\beta$ (kinematic) and $\kappa$ (potential). All known GR critical surfaces (photon sphere, ISCO, horizons) emerge as simple fractions of $(\kappa,\beta)$ from the single closure law $\kappa^2=2\beta^2$ (topologically derived virial-like theorem, Theorem~\ref{thm:Closure}).

\textbf{\textit{WILL Part~I} offers new perspective on several long-standing problems, including:}
\begin{itemize}[nosep,leftmargin=*]
    \item \textbf{Resolution of the classical $M\sin i$ degeneracy} (by introducing independent invariants, \S\ref{sec:M_sin(i)}),
    \item \textbf{Resolution of GR singularities} (via naturally bounded $\rho_{\max}=c^{2}/(8\pi G r^{2})$, \S\ref{sec:no_singularities}),
    \item \textbf{Derivation of the equivalence principle} (from common rest-invariant scaling, \S\ref{thm:equivalence}),
    \item \textbf{Removal of local energy density ambiguity} $\rho=\kappa^{2}c^{2}/(8\pi G r^{2})$ (\S\ref{lem:norm_id}),
    \item \textbf{Clear relational symmetry between kinematic and potential} projections (\S\ref{projections_symmetry}),
    \item \textbf{A closed system of equations for Relational Orbital Mechanics (R.O.M.)} without $G$, mass, or differential formalism (\S\ref{sec:rom}),
    \item \textbf{A computationally simpler and ontologically consistent foundation} for subsequent papers on cosmology (\href{https://willrg.com/documents/WILL_RG_II.pdf}{Part~II}) and quantum mechanics (\href{https://willrg.com/documents/WILL_RG_III.pdf}{Part~III}).
\end{itemize}
\end{minipage}
\end{center}

\vspace{12pt}
\hrule
\vspace{8pt}

% --- Begin two-column body ---
\begin{multicols}{2}

% =============================================================================
% SECTION 1: FOUNDATIONAL PRINCIPLES
% =============================================================================

\begin{quote}
\textit{``There is no such thing as an empty space, i.e., a space without field.\ \ldots\ Space-time does not claim existence on its own, but only as a structural quality of the field.''}
\end{quote}
\noindent\textit{--- Albert Einstein, \textbf{Relativity: The Special and the General Theory} (Appendix~V), 1952 edition, p.~155.}

\begin{tcolorbox}[colback=gray!5, colframe=black!80!black, title=Reading Protocol]
This document must be read \textbf{literally}. All terms are defined within the relational framework of WILL Relational Geometry. Any attempt to reinterpret them through conventional notions (\emph{absolute energies, external backgrounds, hidden containers}) will produce distortions and misreadings.
\end{tcolorbox}

\section{Foundational Principles}
\begin{quote}
  \textit{This Approach Does not Describe Physics; it Generates it.}
\end{quote}

\begin{tcolorbox}[colback=gray!5, colframe=black!80!black, title=Guiding Principle:]
\[
\boxed{\textbf{Nothing is assumed. Everything is derived.}}
\]
\textbf{\textit{\href{https://willrg.com/LOGOS_MAP/}{Interactive Derivation Map}}}
\end{tcolorbox}

\begin{principle}[Epistemic Hygiene]
\hypertarget{pr:epistemic}{}\label{pr:epistemic}
\emph{Epistemic Hygiene as Refusal to Import Unjustified Assumptions.} This line of reasoning derives physics by \textbf{removing hidden assumptions}, rather than introducing new postulates. This construction is deliberate and contains zero free parameters. No assumptions are introduced and no constructs are retained unless they are geometrically or energetically necessary.
\end{principle}

\begin{principle}[Ontological Minimalism]
\hypertarget{pr:minimalism}{}\label{pr:minimalism}
Any fundamental theory must proceed from the minimum possible number of ontological assumptions. The burden of proof lies with any assertion that introduces additional complexity or new entities.
\end{principle}

\subsection*{No Ontological Commitments}
This model makes no ontological claims about the ``existence'' of space, particles, or fields. All phenomena are treated as observer-dependent relational projections.

\begin{principle}[Relational Origin]
\hypertarget{pr:relational}{}\label{pr:relational}
\emph{All physical quantities must be defined by their relations.} Any introduction of absolute properties risks reintroducing metaphysical artefacts and contradicts the foundational insight of relationalism.
\end{principle}

\begin{principle}[Simplicity]
\hypertarget{pr:simplicity}{}\label{pr:simplicity }
\emph{\textbf{Everything} must be expressed in the \textbf{simplest} form possible.} Any unjustified complexity risks reintroducing metaphysical artefacts.
\end{principle}

\begin{principle}[Mathematical Transparency]
\hypertarget{pr:Mathematical}{}\label{pr:Mathematical}
\begin{enumerate}[nosep]
    \item Every mathematical phrase carries its own ontological statement.
    \item Each mathematical object must correspond to an explicitly identifiable relation between observers.
    \item Every symbol must be anchored to a unique physical idea.
    \item Introducing symbols without explicit necessity constitutes \textbf{over-parameterization}.
    \item Number of symbols = Number of independent physical ideas.
\end{enumerate}
\end{principle}


% =============================================================================
% SECTION 2: WHAT IS ENERGY?
% =============================================================================

\subsection{What is Energy in the Relational Framework?}

Across all domains of physics, one empirical fact persists: in every closed system there exists a quantity that never disappears or arises spontaneously, but only transforms in form. This invariant is observed under many guises---kinetic, potential, thermal, quantum---yet all are interchangeable, pointing to a single underlying structure. Crucially, this quantity is never observed directly, but only through \textit{differences between states}. Its value is relational, not absolute.

\begin{tcolorbox}[colback=gray!5, colframe=black!80!black, title=Energy:]
\begin{definition}[Energy]
\hypertarget{def:energy}{}\label{def:energy}
\[
\boxed{
\begin{array}{c}
\textbf{Energy is the invariant relational measure} \\
\text{of state transformation} \\
\text{within a topologically closed system.}
\end{array}
}
\]
\end{definition}
\end{tcolorbox}


% =============================================================================
% SECTION 3: SUBSTANTIALISM VS. RELATIONALISM
% =============================================================================

\section{The No-Go Theorem for Fundamental One-Point Dynamics}
\hypertarget{sec:sub_v_rel}{}\label{sec:sub_v_rel}
\hypertarget{thm:no-go}{}\label{thm:no-go}

\begin{definition}[Self-Centering]
Each observer defines itself as the relational origin:
\[
(\beta,\kappa)=(0,0).
\]
This is an ontological definition, not a coordinate choice.
\end{definition}

\begin{definition}[Relational Reciprocity]
The only invariant quantity between two observers is the norm of the Total Relational Shift:
\[
Q^2=\beta^2+\kappa^2.
\]
Reciprocity is invariance of this norm under the self-centering operation performed independently by each observer.
\end{definition}

\begin{definition}[Absence of Background]
There exists no shared background structure: no global state space, no external time parameter, and no common coordinate system simultaneously hosting the states of distinct observers.
\end{definition}

\begin{definition}[Operationality]
A physical quantity is admissible only if it is either (1)~directly measurable, or (2)~algebraically derivable from measurable quantities, without invoking non-observable auxiliary structures.
\end{definition}

\begin{definition}[One-Point Dynamics]
By one-point dynamics we mean any formulation in which a system is represented by a state $x$ in a global space, physical law is given by a local evolution rule $\dot{x}=F(x)$, and temporal evolution is defined as transitions of the same point through neighboring states. This includes Lagrangian, Hamiltonian, field-theoretic, and metric-based formulations.
\end{definition}

\begin{theorem}[No-Go for Fundamental One-Point Dynamics]
Under Self-Centering, Relational Reciprocity, Absence of Background, and Operationality, no one-point dynamical formulation can be simultaneously: (1)~relationally reciprocal, (2)~operationally well-defined, (3)~background-independent, (4)~ontologically minimal. Therefore one-point dynamics cannot be fundamental.
\end{theorem}

\begin{proof}
A one-point evolution law implicitly assumes a persistent identity of the system across states: it must be meaningful to say that the system at $x(t_1)$ and at $x(t_2)$ is \emph{the same system}. This is a substantialist commitment---an intrinsic identity independent of relational context.

Under Absence of Background, no such structure exists: states are defined only through mutual relational shifts $Q$. The continuous trajectory $x(t)$---and with it $\dot{x}$ or $\delta x$---is an auxiliary construct requiring non-observable distinctions between arbitrarily close states, inadmissible by Operationality.

By Self-Centering, an observer is always at $(\beta,\kappa)=(0,0)$ in its own relational origin: there is no operationally available notion of an observer's own worldline as a sequence of distinct states inside a shared arena.

By Relational Reciprocity, only the scalar norm $Q$ is invariant under mutual self-centering: directional quantities required by one-point dynamics---such as signed increments, tangent vectors, gradients, forces, or local generators---are not reciprocity-invariant and therefore cannot be fundamental.

To restore a well-defined one-point dynamics, one must add a global time parameter, a shared state manifold, a persistent identity map, or a background metric. Each addition violates Absence of Background.
\end{proof}

\begin{corollary}[Constraint-Based Fundamental Law]
Under the same assumptions, admissible fundamental laws must be algebraic and relational: they constrain mutual states through reciprocity-invariant quantities rather than prescribing one-point evolution.
\end{corollary}


\subsection{Strict Minimality of Relational Constraint Laws}
\hypertarget{thm:minimality}{}\label{thm:minimality}

\begin{theorem}[Strict Minimality of Relational Constraint Laws]
Among all formulations capable of reproducing the same observable predictions, relational constraint-based laws have strictly lower ontological cost than any one-point dynamical formulation.
\end{theorem}

\begin{proof}
Relational constraint laws achieve closure using only self-centering, reciprocity, and algebraic invariants such as
\href{https://willrg.com/documents/WILL_RG_I.pdf#sec:DisQ}{$Q^2=\beta^2+\kappa^2$},
\href{https://willrg.com/documents/WILL_RG_I.pdf#thm:Closure}{$\kappa^2=2\beta^2$}, and
\href{https://willrg.com/documents/WILL_RG_I.pdf#}{$\Delta E_{A \to B} + \Delta E_{B \to A} = 0$}.
No background structures, generators, or evolution parameters are required.

Any one-point dynamical formulation requires at least one additional ontological primitive (global state space, ordering parameter, identity map, or generator) not present in the relational scheme. Since relational laws reproduce the empirical content of one-point dynamics (as demonstrated in the WILL Trilogy, \ref{trilogy}) with strictly fewer primitives, the relational formulation is minimal.
\end{proof}

\begin{tcolorbox}[colback=gray!5, colframe=black!80!black, title=Minimality Result]
\[
\boxed{
\text{Relational Constraints}
\;<\;
\text{One-Point Dynamics}
}
\]
The inequality denotes strict ontological and operational minimality.
\end{tcolorbox}


\subsection{General Consequence}
\hypertarget{sec:general_consequence}{}\label{sec:general_consequence}

The adoption of substantialist assumptions produces three structural consequences: (1)~\textit{Inflated Formalism}: equations multiply to compensate for ontological error; (2)~\textit{Loss of Transparency}: physical meaning becomes hidden behind coordinate dependencies; (3)~\textit{Empirical Fragmentation}: each domain requires separate constants. By contrast, Relationalism---as epistemic hygiene---enforces relational closure and yields simplicity through necessity.

\[
\boxed{\text{Surplus Ontology}} \Longrightarrow
\boxed{\text{Ontological Duplication}} \Longrightarrow
\boxed{\text{Mathematical Inflation}}
\]


% =============================================================================
% SECTION 4: THE ONTOLOGICAL BLIND SPOT
% =============================================================================

\section{Ontological Blind Spot in Modern Physics}

The standard formulation of General Relativity relies on the concept of a pre-existing, container-like spacetime that then gets ``filled'' with fields and matter. This is in direct tension with Principle~\ref{pr:relational}. The standard derivation begins with the Einstein--Hilbert action built upon the metric as the fundamental variable, defined on a smooth manifold assumed to exist \textit{a priori}. This manifold carries topological and differential structure---an absolute scaffold violating Principle~\ref{pr:minimalism}.

In the standard formulation, the stress-energy tensor is derived from the variation of the matter Lagrangian with respect to the metric. This assumes that energy is a property of matter fields that can be localized in spacetime. However, this localization is frame-dependent and leads to the non-uniqueness of the gravitational energy-momentum pseudotensor---a direct violation of Principle~\ref{pr:relational}.

\subsection{Historical Pattern: Breakthroughs Delete, Not Add}
\begin{itemize}[nosep]
  \item \textbf{Copernicus} eliminated the Earth/cosmos separation.
  \item \textbf{Newton} eliminated the terrestrial/celestial law separation.
  \item \textbf{Einstein} eliminated the space/time separation.
  \item \textbf{Maxwell} eliminated the electricity/magnetism separation.
\end{itemize}
The spacetime--energy split is the only survivor of this pruning sequence.

\subsection{The Contemporary Split}
All present-day theories (SR, GR, QFT, $\Lambda$CDM, Standard Model) are built with a bi-variable syntax:
\[
\underbrace{\text{fixed manifold + metric}}_{\text{structure}}
\;+\;
\underbrace{\text{fields + constants}}_{\text{dynamics}}.
\]
No observation demands this duplication; it is retained purely because the resulting Lagrangians are empirically adequate \emph{inside} the split. The split is not an empirical discovery but an unpaid ontological debt.


% =============================================================================
% SECTION 5: THE UNIFYING PRINCIPLE
% =============================================================================

\section{The Unifying Principle: Removing the Hidden Assumption}

\begin{lemma}[False Separation]
\hypertarget{lem:false-separation}{}\label{lem:false-separation}
Any model that treats processes as unfolding \emph{within} an independent background necessarily assigns to that background structural features not derivable from relations among the processes themselves. Such a background constitutes an \emph{extraneous absolute}.
\end{lemma}

\begin{proof}
Suppose an independent background exists. Then at least one of its structural attributes---metric relations, a preferred orientation, or a class of inertial frames---remains fixed regardless of interprocess data. This attribute is not relationally inferred but posited \textit{a priori}, violating relational closure.
\end{proof}

\begin{corollary}[Structure-Dynamics Coincidence]
\hypertarget{cor:coincidence}{}\label{cor:coincidence}
To avoid the artifact of Lemma~\ref{lem:false-separation}, the structural arena and the dynamical content must be identified: geometry \emph{is} energy, and energy \emph{is} geometry.
\end{corollary}

\begin{principle}[Ontological Principle: Removing the Hidden Assumption]
\hypertarget{pr:unifying}{}\label{pr:unifying}
\[
\boxed{\textbf{SPACETIME} \;\equiv\; \textbf{ENERGY}}
\]
This is not introduced as a new ontological entity but as a Principle with \emph{negative ontological weight}: it removes the hidden unjustified separation between ``structure'' and ``dynamics.''
\end{principle}

\begin{remark}[Auditability]
Principle~\ref{pr:unifying} is foundational but testable: subject to (i)~geometric audit (internal logical consequences) and (ii)~empirical audit (agreement with data).
\end{remark}

\begin{definition}[WILL]
\hypertarget{def:will}{}\label{def:will}
$\textbf{WILL}\equiv\textbf{SPACE-TIME-ENERGY}$ is the technical term for the unified relational structure determined by Principle~\ref{pr:unifying}. All physically meaningful quantities are relational features of WILL; no external container is permitted.
\end{definition}


% =============================================================================
% SECTION 6: DERIVING THE WILL STRUCTURE
% =============================================================================

\section{Deriving the WILL Structure}
\hypertarget{sec:WILL_structure}{}\label{sec:WILL_structure}

Having established Principle~\ref{pr:unifying}, we derive its necessary geometric and physical consequences.

\begin{lemma}[Closure]
\hypertarget{lem:closure}{}\label{lem:closure}
Under Principle~\ref{pr:unifying}, WILL is self-contained: there is no external reservoir into or from which the relational resource can flow.
\end{lemma}

\begin{proof}
If WILL were not self-contained, an external structure would mediate exchange, serving as a background distinct from the dynamics, contradicting Corollary~\ref{cor:coincidence}.
\end{proof}

\begin{lemma}[Conservation]
\hypertarget{lem:conservation}{}\label{lem:conservation}
Within WILL, the total relational ``transformation resource'' (energy) is conserved.
\end{lemma}

\begin{proof}
By Lemma~\ref{lem:closure}, no external fluxes exist. Any change in one part of WILL must be balanced by complementary change elsewhere.
\end{proof}

\begin{lemma}[Isotropy from Background-Free Relationality]
\hypertarget{lem:isotropy}{}\label{lem:isotropy}
If no external background is allowed (Cor.~\ref{cor:coincidence}), then no direction can be \emph{a priori} privileged. Thus the admissible relational geometry must be maximally symmetric.
\end{lemma}

\begin{proof}
Any alleged privilege not constructible from relations among participants is unobservable (pure gauge). Therefore the carrier must have no intrinsic privileged direction or point.
\end{proof}

\subsection{Derivation of the Relational Carriers}

\begin{theorem}[Minimal Relational Carriers of the Conserved Energy Resource]
\hypertarget{thm:carriers}{}\label{thm:carriers}
The minimal relational carriers satisfying Closure, Conservation, and maximal Symmetry (Lemmas~\ref{lem:closure}--\ref{lem:isotropy}) are:
    \begin{enumerate}[label=(\alph*),nosep]
        \item $S^1$ for \emph{directional} (kinematic) relational transformation;
        \item $S^2$ for \emph{omnidirectional} (gravitational) relational transformation.
    \end{enumerate}
\end{theorem}

\begin{proof}
\textbf{(a) Directional (Kinematic) Relation:}
The simplest non-trivial 1-DOF relation: transformation from State~A to State~B. Per Principle~\ref{pr:relational}, this interaction can only be described from the frame of A or B. From B's perspective, any complex 3D motion of A is operationally perceived, within 1-DOF, only as a change in the rate of approach or recession. Thus the fundamental description is necessarily one-dimensional. By Lemma~\ref{lem:closure} this 1D geometry must be closed; by Lemma~\ref{lem:isotropy} it must be maximally symmetric. The unique carrier is $S^1$.

\textbf{(b) Omnidirectional (Gravitational) Relation:}
The other minimal 2-DOF type: a central state relating to the locus of all equidistant states. By Lemma~\ref{lem:isotropy}, the conserved resource must be distributed uniformly across all orientations. The minimal 2D carrier that is closed and maximally symmetric is $S^2$.
\end{proof}

\begin{corollary}[Uniqueness]
\hypertarget{cor:uniqueness}{}\label{cor:uniqueness}
Under Principle~\ref{pr:unifying} with Lemmas~\ref{lem:closure}--\ref{lem:isotropy}, $S^1$ and $S^2$ are \emph{necessary} relational carriers for directional and omnidirectional modes of energy transformation.
\end{corollary}

\begin{remark}[Non-spatial Reading]
\hypertarget{rem:nonspatial}{}\label{rem:nonspatial}
Throughout this paper, $S^1$ and $S^2$ are not spacetime geometries. They are relational carriers encoding closure, conservation, and isotropy. Ordinary spatial and temporal notions are emergent descriptors.
\end{remark}


% =============================================================================
% SECTION 7: AMPLITUDE-PHASE DUALITY
% =============================================================================

\section{The Amplitude-Phase Duality}

\begin{lemma}[Duality of Evolution]
\hypertarget{lem:duality}{}\label{lem:duality}
The identification of spacetime with energy necessitates two complementary relational measures: (1)~the \textbf{amplitude} of transformation (external shift), and (2)~the \textbf{phase} of transformation (internal order).
\end{lemma}

\begin{proof}
Any complete description of transformation must specify both what changes and how that change is internally ordered. Relational Carriers $S^1$ and $S^2$ each provide the minimal geometry enforcing such complementarity: their orthogonal projections furnish precisely two non-redundant coordinates.
\end{proof}

\begin{definition}[The Amplitude Projection (External Interaction)]
Denoted by $\beta$ (kinematic) and $\kappa$ (gravitational). This component measures the \textit{extent of external relation}, manifesting as momentum or potential intensity.
$\text{Amplitude} \to \text{External Power (Kinetic/Potential)}$.
\end{definition}

\begin{definition}[The Phase Projection (Internal Evolution)]
Denoted by $\beta_Y$ and $\kappa_X$. This component measures the \textit{internal ordering}, governing proper time and proper length. Phase~$=1$ represents maximal internal flow (rest/vacuum); Phase~$=0$ represents cessation of internal causality (light-speed/horizon).
$\text{Phase} \to \text{Internal Order (Time/Structure)}$.
\end{definition}

\begin{theorem}[Universal Conservation of Relation]
For both kinematic ($S^1$) and gravitational ($S^2$) modes:
\begin{equation}
\underbrace{\text{Amplitude}^2}_{\text{External}} + \underbrace{\text{Phase}^2}_{\text{Internal}} = 1
\end{equation}
\end{theorem}

\begin{proof}
By the geometric nature of the relational carriers $S^1$ ($\beta^2+\beta_Y^2=1$) and $S^2$ ($\kappa^2+\kappa_X^2=1$), encoding the finite relational budget.
\end{proof}

\begin{proposition}[Physical Interpretation: Relativistic Effects]
\hypertarget{prop:relativity}{}\label{prop:relativity}
The conservation law implies that redistribution between amplitude and phase manifests as time dilation and length contraction.
\end{proposition}

\begin{proof}
An increase in Amplitude ($\beta$, $\kappa$) enforces a decrease in Phase ($\beta_Y$, $\kappa_X$). This reduction corresponds to dilation of proper time and contraction of proper length. The relativistic and gravitational trade-off is the direct physical expression of geometric closure.
\end{proof}


% =============================================================================
% SECTION 8: ENERGY AS A RELATION
% =============================================================================

\section{Energy as a Relation: What $\kappa$ and $\beta$ Actually Mean}
\hypertarget{sec:kappabeta}{}\label{sec:kappabeta}

In the relational framework, physical parameters like energy, speed, and gravitational potential do not belong to objects. They represent how an observer measures state differences from their own point of view. Your perspective is always the reference frame---you are always at the origin $(0,0)$ on your $(\beta,\kappa)$ plane.

$\beta$ measures how much of the universal ``speed of change'' you see as motion through space, relative to yourself. $\kappa$ measures how deeply an object sits in a gravitational field, as seen from your position. Energy is the capacity to move between states; saying ``the object's energy'' implicitly means ``the object's energy as measured from your perspective.''


% =============================================================================
% SECTION 9: KINETIC ENERGY PROJECTION ON S^1
% =============================================================================

\section{Kinetic Energy Projection on $S^1$}

\begin{tcolorbox}[colback=gray!5, colframe=black!80!black, title=Desmos project]
\textbf{\href{https://www.desmos.com/geometry/6k1um2qbzm}{Kinetic projection $\beta$ on the $S^1$ carrier (Special Relativity)}}
\end{tcolorbox}

Since $S^{1}$ encodes one-dimensional shift, the total energy $E$ projects consistently onto both axes:
$E_{X} = E\beta$, $E_{Y} = E\beta_{Y}$.

\subsection{The Geometric Nature of Mass}

\begin{theorem}[Invariant Projection of Rest Energy]
\hypertarget{thm:restenergy}{}\label{thm:restenergy}
For any state $(\beta,\beta_Y)$ on $S^1$, the total energy $E$ must scale such that its vertical projection remains constant and equal to the rest energy $E_0$:
\[
E\beta_Y = E_0.
\]
\end{theorem}

\begin{proof}
Let $E$ be distributed on $S^1$ with projections $\beta$ and $\beta_Y$. By Principle~\ref{pr:relational}, the rest energy $E_0$ is defined by internal relations and must be invariant under changes in external relation. In self-reference, $\beta=0 \Longrightarrow \beta_Y=1$, so:
\[
\boxed{E_{Y} \equiv E_0}
\]
Geometric consistency requires:
$E\beta_Y = E_0 \Longrightarrow E = E_0/\beta_Y$.
The ``hypotenuse'' (Total Energy $E$) is not a fixed-length vector but a scalable relational magnitude that grows to preserve the invariant leg $E_0$.
\end{proof}

The historical Lorentz factor $\gamma$ is the reciprocal of $\beta_Y$: $\gamma = 1/\beta_Y$.

\begin{corollary}[Rest Energy and Mass Equivalence]
\hypertarget{cor:restmass}{}\label{cor:restmass}
Within $c=1$, $E_{0} = m$. Mass is not independent but the rest-energy invariant itself.
\end{corollary}

\subsection{Energy-Momentum Relation}

\begin{tcolorbox}[colback=gray!5, colframe=black!80!black, title=Desmos project]
\textbf{\href{https://www.desmos.com/geometry/mbkzikthkh}{Energy-momentum Triangle}}
\end{tcolorbox}

\begin{proposition}[Horizontal Projection as Momentum]
\hypertarget{prop:momentum}{}\label{prop:momentum}
$p \equiv E\beta$ \quad ($c=1$).
\end{proposition}

\begin{proof}
The rest state is $(\beta,\beta_Y)=(0,1)$. A shift measure must vanish at rest, grow monotonically with $|\beta|$, and flip sign under $\beta\mapsto -\beta$. The only candidate is $E\beta$.
\end{proof}

\begin{corollary}[Energy-Momentum Relation]
\hypertarget{cor:energymomentum}{}\label{cor:energymomentum}
\[
E^{2} = p^{2} + m^{2} \quad (c=1), \qquad
E^{2} = (pc)^{2} + (mc^{2})^{2}.
\]
\end{corollary}

\begin{proof}
By closure, $(E\beta)^2 + (E\beta_Y)^2 = E^2$. Substituting $p=E\beta$ and $m=E_0$ proves the claim.
\end{proof}

\begin{table}[h]
\centering\footnotesize
\begin{tabular}{|c|c|}\hline
\multicolumn{2}{|c|}{$\beta=v/c$, \quad $\theta_1= \arccos(\beta)$}\\
\hline
\textbf{Algebraic Form} & \textbf{Trigonometric Form} \\
\hline
$\beta= v/c = \sqrt{1-\beta_Y^{2}}$& $\beta = \cos(\theta_1)$\\
\hline
$\beta_Y = \sqrt{1-\beta^2}$& $\beta_Y = \sin(\theta_1)$\\\hline
\end{tabular}
\caption{Geometric representation of relativistic effects.}
\end{table}


% =============================================================================
% SECTION 10: POTENTIAL ENERGY PROJECTION ON S^2
% =============================================================================

\section{Potential Energy Projection on $S^2$}

\begin{tcolorbox}[colback=gray!5, colframe=black!80!black, title=Desmos project]
\textbf{\href{https://www.desmos.com/geometry/lsjhnc2e9x}{Potential projection $\kappa$ on $S^2$ carrier}}
\end{tcolorbox}

Analogous to $S^1$, the relational geometry of $S^2$ provides orthogonal projections for omnidirectional transformation:

The \textbf{Amplitude Component} ($\kappa$) represents the relational gravitational measure. $\kappa=1$ denotes saturation (the relational horizon). The \textbf{Phase Component} ($\kappa_X$) governs intrinsic proper length and proper time. They satisfy:
$\kappa_X^2 + \kappa^2 = 1$.

\subsection{Gravitational Meridional Section of $S^2$}
By isotropy the omnidirectional carrier is $S^2$, but any radially symmetric exchange reduces to a great-circle meridional section: $(\kappa_X,\kappa)=(\cos\theta_2,\sin\theta_2)$.

\subsection{Gravitational Tangent Formulation}

In full symmetry with the kinematic case, from $\kappa = \sin\theta_{2}$:
\begin{equation}
    E_{g} = \frac{E_{0}}{\kappa_{X}}, \qquad
    p_g = \frac{E_{0}}{c}\tan\theta_{2},
\end{equation}
yielding $E_{g}^{2} = (p_gc)^{2} + (mc^{2})^{2}$.


% =============================================================================
% SECTION 11: GEOMETRIC COMPOSITION + TOTAL SHIFT Q
% =============================================================================

\section{Geometric Composition of SR and GR Factors}
\hypertarget{sec:geometric_composition}{}\label{sec:geometric_composition}

On the unit kinematic circle ($S^1$): $(\beta,\beta_Y)=(\cos\theta_1,\sin\theta_1)$, giving
\[
E=\frac{E_0}{\beta_Y}=\frac{E_0}{\sin\theta_1}, \qquad
p=E_0\cot\theta_1.
\]

On the gravitational circle ($S^2$): $(\kappa_X,\kappa)=(\cos\theta_2,\sin\theta_2)$, giving
\[
E_g=\frac{E_0}{\kappa_X}=\frac{E_0}{\cos\theta_2}, \qquad
p_g=E_0\tan\theta_2.
\]

\subsection{Total Relational Shift $Q$}
\hypertarget{sec:DisQ}{}\label{sec:DisQ}

\begin{tcolorbox}[colback=gray!5, colframe=black!80!black, title=Desmos project]
\textbf{\href{https://www.desmos.com/geometry/2nkjtezjpi}{$Q$ as Total Relational Shift}}
\end{tcolorbox}

\hypertarget{eq:Q-definition}{}\label{eq:Q-definition}
When an observer observes another system, they assign a Total Relational Shift norm:
\begin{equation}
\boxed{Q^2=\beta^2+\kappa^2}
\end{equation}

Each observer places itself at the relational origin $(\beta,\kappa)=(0,0)$. The other system, looking back, self-centres and applies the same rule, obtaining the same $Q^2$. Thus $Q$ is the norm of the shift, not a spatial distance.

\begin{remark}[Closure-specific simplification]
Under energetic closure $\kappa^{2}=2\beta^{2}$, the norm reduces to $Q^{2}=3\beta^{2}$ (Theorem~\ref{thm:Closure}).
\end{remark}


% =============================================================================
% SECTION 12: RELATIONAL SYMMETRY TABLE
% =============================================================================

\section{Relational Symmetry Between Kinematic and Potential Projections}
\hypertarget{projections_symmetry}{}\label{projections_symmetry}

\begin{table}[h]
\centering\footnotesize
\begin{tabular}{|c|c|}\hline
\multicolumn{2}{|c|}{$\theta_1= \arccos(\beta)$, $\theta_2 = \arcsin(\kappa)$, $\kappa^2 = 2\beta^2$}\\
\hline
\textbf{Algebraic Form} & \textbf{Trigonometric Form} \\\hline
$\beta=v/c$&$\beta=\cos(\theta_1)$\\\hline
$\kappa=\sqrt{R_s/r}$&$\kappa=\sin(\theta_2)$\\\hline
$\beta_Y = \sqrt{1-\beta^2}$& $\beta_Y = \sin(\theta_1)$\\\hline
$\kappa_X = \sqrt{1-\kappa^2}$&$\kappa_X = \cos(\theta_2)$\\\hline
$p=E_0/c \cdot \beta/\beta_Y$&$p=E_0/c \cdot \cot(\theta_1)$\\\hline
$p_g=E_0/c \cdot \kappa/\kappa_X$&$p_g=E_0/c \cdot \tan(\theta_2)$\\\hline
$\tau = \beta_Y\kappa_X$&$\tau = \sin(\theta_1) \cos(\theta_2)$\\\hline
$Q =\sqrt{\kappa^2 + \beta^2}=\sqrt{3}\beta$&$Q = \sqrt{3} \cos(\theta_1)$\\\hline
\end{tabular}
\caption{Unified representation of relativistic and gravitational effects for closed systems.}
\end{table}


% =============================================================================
% SECTION 13: CRITICAL RADII
% =============================================================================

\begin{table}[h]
\centering\footnotesize
\renewcommand{\arraystretch}{1.1}
\begin{tabular}{|l|c|c|c|c|}
\hline
\textbf{Phenomenon} & $r$ & $\beta^2$ & $\kappa^2$ & $Q^2$ \\
\hline
ISCO (Static) & $3R_s$ & $0$ & $1/3$ & $1/3$ \\
ISCO (Dynamic) & $3R_s$ & $1/6$ & $1/3$ & $1/2$ \\
\hline
Photon sphere (Static) & $\frac{3}{2}R_s$ & $0$ & $2/3$ & $2/3$ \\
Causal Dynamic horizon & $\frac{3}{2}R_s$ & $1/3$ & $2/3$ & $1$ \\
\hline
Causal Static horizon & $R_s$ & $0$ & $1$ & $1$ \\
Dynamic horizon & $R_s$ & $1/2$ & $1$ & $3/2$ \\
\hline
Extremal Kerr horizon & $R_s/2$ & $1$ & $2$ & $3$ \\
\hline
\end{tabular}
\caption{Critical radii and their projectional parameters.}
\end{table}


\subsection{Kinematic Mechanism of the Photon Sphere}
\hypertarget{sec:photon_sphere}{}\label{sec:photon_sphere}

At the Causal Circular Horizon ($Q^2=1$), with closure $\kappa^2 = 2\beta^2$: $\beta^2 = 1/3$, $\kappa^2 = 2/3$. For circular orbit ($e = 0$):
\[
\Delta\varphi = \frac{2\pi \cdot 3 (1/3)}{1 - 0} = 2\pi.
\]
An orbital precession of exactly $2\pi$ per revolution means the trajectory continuously folds onto itself, creating a closed spherical locus---the photon sphere. This is the true kinematic definition: not the speed of the particle, but the unitary topological closure of its trajectory.


% =============================================================================
% SECTION 14: ENERGETIC CLOSURE CONDITION
% =============================================================================

\section{Unification of Projections: The Energetic Closure Condition}

\begin{remark}[From slice to whole $S^2$]
Although we parametrise a single meridional great circle, $\kappa^2$ denotes the total omnidirectional budget. The factor~$2$ reflects that $S^2$ has two independent DOF.
\end{remark}

\begin{lemma}[DOF-Indifference]
\hypertarget{lem:dof-indifference}{}\label{lem:dof-indifference}
Under maximal symmetry and ontological minimalism, any admissible conserved budget must assign equal quadratic weight to each independent DOF.
\end{lemma}

\begin{proof}
If two independent DOF contributed unequal weights, the theory would contain an implicit weighting structure that distinguishes DOF---a privileged feature violating maximal symmetry.
\end{proof}

\begin{theorem}[Closure]
\hypertarget{thm:Closure}{}\label{thm:Closure}
The only exchange rate between $S^1$ (1~DOF) and $S^2$ (2~DOF) compatible with closure, maximal symmetry, and minimalism is:
\[
\boxed{\kappa^2 = 2\beta^2.}
\]
\end{theorem}

\begin{proof}
Let $b$ denote the conserved quadratic budget per DOF (Lemma~\ref{lem:dof-indifference}). Then $B_{S^1}=b$ and $B_{S^2}=2b$. The exchange rate $\mathcal{R}=B_{S^2}/B_{S^1}=2$, hence $\kappa^2=2\beta^2$. Any $\mathcal{R}\neq 2$ would require hidden structure.
\end{proof}

\hypertarget{def:closure_factor}{}\label{def:closure_factor}
\begin{definition}[Closure Factor]
$\delta \equiv \kappa^2/(2\beta^2)$.
A subsystem is energetically closed if $\langle\delta\rangle_{\text{cycle}}=1$. For circular orbits, $\delta \equiv 1$.
\end{definition}

\begin{corollary}[Energetic Closure Criterion]
\hypertarget{cor:closurecriterion}{}\label{cor:closurecriterion}
Closed systems satisfy $\kappa^2 = 2\beta^2$ identically. Open systems display $\delta \neq 1$; when all channels are included, closure is restored.
\end{corollary}


% =============================================================================
% SECTION 15: ENERGY-SYMMETRY LAW
% =============================================================================

\section{Energy-Symmetry Law}
\hypertarget{sec:energy-symmetry}{}\label{sec:energy-symmetry}

\begin{theorem}[Energy Symmetry]
\hypertarget{thm:conservation}{}\label{thm:conservation}
The specific energy differences perceived by two observers for a transition between their states balance:
\begin{equation}
\Delta E_{A \to B} + \Delta E_{B \to A} = 0.
\end{equation}
\end{theorem}

\begin{proof}
Consider observer~$A$ at rest at radius $r_A$ and observer~$B$ orbiting at radius $r_B$ with velocity $v_B$.

From $A$'s perspective:
\begin{equation}
\boxed{\Delta E_{A \to B} = \frac{1}{2}(\kappa_A^2 - \kappa_B^2) + \frac{1}{2}(\beta_B^2 - \beta_A^2)}
\end{equation}

From $B$'s perspective:
\[
\Delta E_{B \to A} = \frac{1}{2}((\kappa_B^2 - \kappa_A^2) + (\beta_A^2 - \beta_B^2))
\]

Summing: $\Delta E_{A \to B} + \Delta E_{B \to A} = 0$.
\end{proof}

\subsection{The Specific Energy Transfer}
The energy transfer is:
\[
\Delta E_{A \to B} = \frac{1}{2}(\kappa_A^2 - \kappa_B^2) + \frac{1}{2}(\beta_B^2 - \beta_A^2).
\]

Under closure ($\kappa^2=2\beta^2$), simplified cases emerge:

\textbf{Surface-to-Orbit}: $E_{A\to B}/E_{0B} = \frac{1}{2}(\kappa_A^2 - \beta_B^2)$.

\textbf{Orbit-to-Orbit}: $E_{A\to B}/E_{0B} = \frac{1}{2}(\beta_A^2 - \beta_B^2)$.

\subsection{Universal Speed Limit}
\begin{theorem}[Universal Speed Limit]
The condition $\beta \leq 1$ ($v \leq c$) is an intrinsic requirement for maintaining causal and energetic consistency.
\end{theorem}

\begin{proof}
If $\beta > 1$, the kinetic budget $\frac{1}{2}\beta^2$ becomes arbitrarily large. No finite process could provide a balancing reverse transfer. The Energy-Symmetry Law would be violated.
\end{proof}

\subsection{Single-Axis Energy Transfer and the Nature of Light}
\hypertarget{sec:nature_of_light}{}\label{sec:nature_of_light}

\begin{theorem}[Single-Axis Transformation Principle]
\hypertarget{thm:single-axis}{}\label{thm:single-axis}
For light: $\beta = 1 \Longrightarrow \beta_Y = 0$.
All transformation occurs along a single axis.
\end{theorem}

\begin{proof}
At $\beta=1$, the Phase component vanishes ($\beta_Y=0$): no rest frame exists, eliminating the dual-framing that justifies the $\frac{1}{2}$ partitioning. The energy invariant for photons becomes $W_\gamma = \kappa^2 - 1$ (not $\frac{1}{2}(\kappa^2-1)$). Consequently:
\[
\Phi_\gamma = \kappa^2 c^2, \qquad \Phi_{\text{mass}} = \tfrac{1}{2}\kappa^2 c^2.
\]
This factor-of-2 explains the experimentally verified doubled gravitational effect on light.
\end{proof}


% =============================================================================
% SECTION 16: EQUIVALENCE PRINCIPLE
% =============================================================================

\section{Equivalence Principle as Derived Identity}

\begin{lemma}[Unified Relational Scaling]
\hypertarget{lem:unifiedscaling}{}\label{lem:unifiedscaling}
Both transformations act as independent projections of the same invariant $E_0$:
$E = E_0/\beta_Y$, \quad $E_g = E_0/\kappa_X$.
\end{lemma}

\begin{theorem}[Equivalence of Inertial and Gravitational Response]
\hypertarget{thm:equivalence}{}\label{thm:equivalence}
Composing the independent stretches yields:
\[
E_{\mathrm{loc}} = \frac{E_0}{\beta_Y\,\kappa_X} = \frac{E_0}{\sqrt{(1-\beta^2)(1-\kappa^2)}},
\]
with effective mass $m_{\mathrm{eff}} = E_0/(\beta_Y\kappa_X c^2)$. Therefore:
\[
\boxed{m_g \equiv m_i \equiv m_{\mathrm{eff}}},
\]
and the Einstein equivalence principle follows as a structural identity of WILL.
\end{theorem}

\begin{corollary}[Mass Invariance under Relational Scaling]
\hypertarget{cor:massinvariance}{}\label{cor:massinvariance}
\[
\boxed{m_g \equiv m_i \equiv m = E_0/c^2}
\]
is not a dynamical statement but the definition of rest invariance.
\end{corollary}


% =============================================================================
% SECTION 17: CLASSICAL LIMITS
% =============================================================================

\section{Classical Limits of the Energy-Symmetry Law}
\hypertarget{sec:classical_vs_WILL}{}\label{sec:classical_vs_WILL}

\subsection{Keplerian Energy and Minkowski-like Form}

For a test body on a circular orbit at radius $a$, defining $\kappa_\oplus^2 \equiv 2GM_\oplus/(R_\oplus c^2)$ and $\beta_{\text{orbit}}^2 \equiv GM_\oplus/(ac^2)$:
\hypertarget{eq:will_minkowski_energy}{}\label{eq:will_minkowski_energy}
\begin{equation}
\frac{E_{\text{tot}}}{E_0} = \frac{1}{2}(\kappa_\oplus^2 - \beta_{\text{orbit}}^2),
\end{equation}
which is structurally identical to a Minkowski interval in $(1+1)$ dimensions.

\subsection{Lagrangian and Hamiltonian as Single-Point Limits}
\hypertarget{sec:LH}{}\label{sec:LH}

The relational Lagrangian represents the kinetic budget at~$B$ relative to the potential budget at~$A$:
\[
\frac{L_{\text{rel}}}{E_0} = \frac{1}{2}(\beta_B^2 + \kappa_A^2).
\]

If one collapses $r_A = r_B = r$, this degenerates into the standard Newtonian Lagrangian $L = \frac{1}{2}mv^2 + GMm/r$. The Hamiltonian follows via Legendre transformation:
\begin{align}
L &\longleftrightarrow \tfrac{1}{2}mc^2(\beta^2 + \kappa^2), \\
H &\longleftrightarrow \tfrac{1}{2}mc^2(\beta^2 - \kappa^2).
\end{align}

The apparent sign difference between $L$ and $H$ is an artifact of this single-point collapse.

\subsection{Newton's Third Law}
\hypertarget{thm:third_law}{}\label{thm:third_law}

\begin{theorem}[Newton's Third Law as a Single-Point Limit]
The Energy-Symmetry Law mathematically necessitates $\vec{F}_{AB} = -\vec{F}_{BA}$ in the classical limit where the two-point relational energy budget is collapsed into $U(\vec{r})$.
\end{theorem}

\begin{proof}
In the collapsed formalism with $U = U(\vec{r}_B - \vec{r}_A)$: $\vec{F}_{AB} = -\nabla_B U = -\nabla U(\vec{r})$ and $\vec{F}_{BA} = -\nabla_A U = +\nabla U(\vec{r})$. Therefore $\vec{F}_{AB} = -\vec{F}_{BA}$.
\end{proof}


% =============================================================================
% SECTION 18: DENSITY, PRESSURE, FIELD EQUATION
% =============================================================================

\section{Derivation of Density, Mass, and Pressure}

\subsection{Derivation of Density}
\hypertarget{sec:density}{}\label{sec:density}

From the projective analysis: $\kappa^2 = R_s/r$, with $R_s = 2Gm_0/c^2$. Translating the 2D parameter $\kappa^2$ on $S^2$ into the conventional 3D volumetric density via the geometric normalization factor $1/(4\pi)$:
\[
\rho = \frac{\kappa^2c^2}{8\pi G r^2}.
\]
At $\kappa^2 = 1$ (horizon): $\rho_{\max} = c^2/(8\pi G r^2)$.
\hypertarget{lem:norm_id}{}\label{lem:norm_id}
\[
\boxed{\kappa^2 = \frac{\rho}{\rho_{\max}} \equiv \Omega}
\]

Self-consistency ($m_0 = \alpha r^n \rho$) requires $n=3$ and $\alpha=4\pi$, yielding $m_0 = 4\pi r^3\rho$.

\subsection{Pressure as Surface Curvature Gradient}
\hypertarget{sec:pressure}{}\label{sec:pressure}

The radial balance relation gives:
\[
P(r) = \frac{c^4}{8\pi G}\frac{1}{r}\frac{d\kappa^2}{dr}.
\]
Using $\kappa^2 = R_s/r$, we find $d\kappa^2/dr = -\kappa^2/r$, hence:
\[
\boxed{P(r) = -\rho(r)\,c^2.}
\]
This negative pressure is not exotic but the geometric tension required for self-consistency when $\kappa^2$ varies radially.


\section{Unified Geometric Field Equation}
\hypertarget{eq:unified_field}{}\label{eq:unified_field}

\[
\boxed{\kappa^2 = \frac{R_s}{r} = \frac{\rho_{\text{field}}}{\rho_{\max}}}
\]

For a static, spherically symmetric configuration with matter density $\rho_{\text{matter}}(r)$:
\hypertarget{eq:will_field_diff}{}\label{eq:will_field_diff}
\begin{equation}
\boxed{\frac{d}{dr}(r\kappa^{2}) = \frac{8\pi G}{c^{2}} r^{2}\rho_{\text{matter}}(r)}
\end{equation}
This reproduces the $tt$-component of the Einstein field equations.

In vacuum ($\rho_{\text{matter}} = 0$): $d(r\kappa^2)/dr = 0 \Longrightarrow r\kappa^2 = R_s$.


\section{No Singularities, No Hidden Regions}
\hypertarget{sec:no_singularities}{}\label{sec:no_singularities}

The constraint $\kappa^2\le 1$ enforces $\rho\le\rho_{\max}$ and $|P|\le c^4/(8\pi G r^2)$, avoiding singularities without extra hypotheses. WILL resolves the singularity problem not by regularizing divergent terms nor by introducing quantum effects, but by geometrically constraining the domain of valid projections. Black holes become energetically saturated but non-singular regions.


\section{Theoretical Ouroboros}
\hypertarget{sec:ouroboros}{}\label{sec:ouroboros}

\hypertarget{eq:unified_identity}{}\label{eq:unified_identity}
We started with $\text{SPACETIME} \equiv \text{ENERGY}$, derived geometry and physical laws, and these laws loop back to define and limit energy and spacetime:
\[
\boxed{\kappa^2 = \frac{R_s}{r} = \frac{\rho}{\rho_{\max}}}
\]
\[
\boxed{
\text{SPACETIME GEOMETRY} \;\equiv\; \text{ENERGY DISTRIBUTION}
}
\]


% =============================================================================
% SECTION 19: R.O.M. --- THE FULL EQUATION SYSTEM
% =============================================================================

\section{Relational Orbital Mechanics (R.O.M.)}
\hypertarget{sec:rom}{}\label{sec:rom}

\begin{tcolorbox}[colback=gray!5, colframe=black!80!black, title=Desmos Project]
\textbf{\href{https://www.desmos.com/calculator/n4lmkpsebx}{R.O.M.}}
\end{tcolorbox}
\hypertarget{eq:rom}{}\label{eq:rom}

\begin{remark}
R.O.M.\ does not describe how a body moves under forces; it classifies the algebraically allowed relational states of a bound system. It is not equations of motion but an algebraically closed system of allowed states within WILL.
\end{remark}

\textbf{Orbital dynamics requires no mass, no $G$, no metric, and no spacetime geometry.} All observable orbital structure follows from two directly measurable frequency projections: $\kappa$ (gravitational projection from redshift) and $\beta$ (kinematic projection from Doppler). Everything else is algebra.

\subsection{Foundational Definitions}

$\boxed{\kappa^2 = 1-\frac{1}{(1+z_{\kappa})^2}}$ \quad ($z_{\kappa}=$ gravitational redshift) \\[3pt]
$\boxed{\beta^2=1-\frac{1}{(1+z_{\beta})^{2}}}$ \quad ($z_{\beta}=$ transverse Doppler shift)

\subsection{Observational $Z$ Inputs}

$Z_{sys}(o)=(1+z_{\kappa o})(1+z_{\beta o})=\tau_{Wo}(o)^{-1}$

$\tau_{Wo}(o)=\kappa_{Xo}(o)\cdot\beta_{Yo}(o)=(Z_{sys}(o))^{-1}$

$z_{\kappa} = 1/\kappa_{X}-1$, \quad $z_{\beta}=1/\beta_{Y}-1$

\subsection{Global System Parameters (Fixed for the Orbit)}

\begin{small}
$\kappa = \sqrt{R_{s}/a} = \sqrt{4W}=\sqrt{1-(1+z_{\kappa})^{-2}}$

$\kappa_{X} = \cos(\theta_2) = \sqrt{1 - \kappa^2}$

$\beta = \kappa/\sqrt{2} = \sqrt{2W}=\sqrt{1-(1+z_{\beta})^{-2}}$

$\beta_{Y} = \sin(\theta_1) = \sqrt{1 - \beta^2}$

$\beta_{\text{int}}=\beta/\sqrt{1-e^2}$

$\tau_{W} = \kappa_{X}\beta_{Y}$

$Q=\sqrt{\kappa^2+\beta^2}=\sqrt{3/2}\,\kappa=\sqrt{3}\,\beta$

$R_{s} = \kappa_{o}^{2} r_o = 2Gm_{0}/c^{2}$

$a = R_{s}/\kappa^{2}$, \quad $m_{0} = \kappa^{2}c^{2}a/(2G) = 4\pi\rho a^{3}$

$t = a/c$, \quad $W =\beta^2/2$

$\Delta\phi = \frac{3\pi}{2}\frac{\kappa_{p}^{4}}{\beta_{p}^{2}} =\frac{2\pi Q^{2}}{1-e^{2}}=6\pi\beta_{\text{int}}^2$

$h_{W} = a \cdot \beta c \cdot e_{Y}$, \quad $\omega = \beta c/a$, \quad $T = 2\pi/\omega$
\end{small}

\subsection{Eccentricity Relations}

$e = 1/\delta - 1 = 2\beta_{p}^{2}/\kappa_{p}^{2}-1$

$e_{Y} = \sqrt{1-e^{2}}$

$e_{X} = (1+e)/(1-e)=r_{a}/r_{p}=\beta_{p}/\beta_{a}=\kappa_{p}^{2}/\kappa_{a}^{2}$

\subsection{Time Integration}

$\omega_{o}(o)=\frac{\beta c}{a}\cdot\frac{(1+e\cos(o))^{2}}{(1-e^{2})^{3/2}}$

$\Delta_{to}(o) = \frac{a}{\beta c}\tau_{o}= \frac{T_{O}}{2\pi}\tau_{o}$

$\tau_{o}=(1-e^{2})^{3/2}\int_{0}^{o}(1+e\cos(\theta))^{-2}d\theta$

\subsection{Perihelion Relations}

$r_{p} = a(1-e)=R_s/\kappa_p^2$

$\kappa_{p} = \kappa\sqrt{1/(1-e)}$, \quad $\kappa_{Xp} = \sqrt{1-\kappa_{p}^2}$

$\beta_{p} = \sqrt{\kappa_{p}^{2}(1+e)/2}$

$\delta_p = 1/(1+e)$, \quad $Q_{p} = \sqrt{\kappa_{p}^{2}+\beta_{p}^{2}}$

\subsection{Aphelion Relations}

$r_a=a(1+e)=R_s/\kappa_a^2$

$\beta_{a} = \beta\sqrt{e_{X}^{-1}}$, \quad $\kappa_{a} = \sqrt{2W+\beta_{a}^{2}}$

$\delta_a = 1/(1-e)$, \quad $Q_{a} = \sqrt{\kappa_{a}^{2}+\beta_{a}^{2}}$

\subsection{Phase Variables (functions of $o$)}

$r_o(o)= a(1-e^{2})/(1+e\cos o) = R_s/\kappa_o^2$

$\kappa_{o} = \sqrt{R_{s}/r} = \kappa_{p}\sqrt{(1+e\cos o)/(1+e)}$

$\kappa_{Xo} = \sqrt{1 - \kappa_o^{2}}$

$\beta_{o}(o) = \sqrt{\kappa_{o}^{2}-2W}$

$\beta_{R}(o)=\beta\,e\sin(o)/\sqrt{1-e^{2}}$

$\beta_{T}(o)=\beta(1+e\cos(o))/\sqrt{1-e^{2}}$

$\beta_{Yo} = \sqrt{1 - \beta_o^{2}}$

$\delta_{o} = (1+e\cos o)/(1+e^{2}+2e\cos o) = \kappa_{o}^2/(2\beta_{o}^2)$

$Q_{o} = \sqrt{\kappa_{o}^{2}+\beta_{o}^{2}}$

$\tau_{Wo}(o)=\kappa_{Xo}\cdot\beta_{Yo}= Z_{sys}(o)^{-1}$

$\Delta_o=3\beta_{\text{int}}^{2} \cdot o$ (precession at phase $o$)

\subsection{Relational Geometry (WILL)}

$\theta_{1} = \arccos(\beta)$, \quad $\theta_{2} = \arcsin(\kappa)$ \quad (non-physical)

$\Delta_{Q} = Q_{o}^{2}-Q^{2}$

$O_{o} = \arccos(-e)$ (balance point where $\kappa_o^2=2\beta_o^2$)

\subsection{Observer-Dependent Relations}

$Z_{raw}(o)=(1+\beta_{\text{int}}(\cos(o+\omega_{i})+e\cos(\omega_{i}))\sin(i))\,Z_{sys}(o)$

$\beta_{los}(o)=\frac{\beta}{\sqrt{1-e^{2}}}(\cos(o+\omega_{i})+e\cos(\omega_{i}))\sin(i)$

$K_{i}=\beta_{\text{int}}\sin(i)$

$Z_{rawmax}=Z_{sys}(o_{i1})(1+K_{i}(1+e\cos\omega_{i}))$

$Z_{rawmin}=Z_{sys}(o_{i2})(1+K_{i}(-1+e\cos\omega_{i}))$

\textbf{Beautiful Identity:}
$\frac{r_{a}}{r_{p}}=\frac{1+e}{1-e}=\frac{\beta_{p}}{\beta_{a}}=\frac{\kappa_{p}^{2}}{\kappa_{a}^{2}}$


% =============================================================================
% SECTION 20: OPERATIONAL INDEPENDENCE FROM G
% =============================================================================

\section{Operational Independence and the Role of Constants}
\hypertarget{sec:operational}{}\label{sec:operational}

\begin{theorem}[Operational Measurability]
\hypertarget{thm:operational}{}\label{thm:operational}
The relational projections are encoded directly in the combined phase interactions of light and are operationally independent of $G$, $c$, or $m_0$.
\end{theorem}

\begin{proof}
\textbf{Step~1.} Spectroscopy measures the total accumulated phase difference. The \textbf{Relational Spacetime Factor} $\tau_W$ is:
\hypertarget{eq:tau_observable}{}\label{eq:tau_observable}
\begin{equation}
\tau_W \equiv \frac{1}{Z_{sys}} = \frac{1}{(1+z_{\kappa})(1+z_{\beta})}
= \underbrace{\kappa_X}_{\text{Grav.\ Phase}} \cdot \underbrace{\beta_Y}_{\text{Kin.\ Phase}}.
\end{equation}

\textbf{Step~2.} The exact identity:
\begin{equation}
\tau_W^2 = 1 - Q^2 + \kappa^2\beta^2,
\end{equation}
demonstrates that the optical signal contains the complete system information, measurable without prior knowledge of mass.
\end{proof}

\subsection{Three Methods for Determining System Scale}

\subsubsection{Method~A: Two-Point Method}

\begin{theorem}[Two-Point Schwarzschild Scale]
\hypertarget{thm:two_point}{}\label{thm:two_point}
\hypertarget{eq:Rs_two_point}{}\label{eq:Rs_two_point}
\begin{equation}
\boxed{R_s = \frac{r_1 r_2}{r_2 - r_1}(\beta_1^2 - \beta_2^2).}
\end{equation}
\end{theorem}

\begin{proof}
By conservation of $W = \frac{1}{2}(\kappa^2-\beta^2)$: $\kappa_1^2 - \kappa_2^2 = \beta_1^2 - \beta_2^2$. Substituting $\kappa^2 = R_s/r$ and solving yields the result.
\end{proof}

\subsubsection{Method~B: Balance Point Method}

\begin{theorem}[Balance Point Formula]
\hypertarget{thm:balance_point}{}\label{thm:balance_point}
\hypertarget{eq:Rs_balance}{}\label{eq:Rs_balance}
\begin{equation}
\boxed{R_s = \frac{a}{2}(3 - \sqrt{1 + 8\tau_W^2(O_o)}).}
\end{equation}
\end{theorem}

\begin{proof}
At $r=a$: $\kappa^2 = R_s/a$, $\beta^2 = R_s/(2a)$. Then $\tau_W^2 = (1-R_s/a)(1-R_s/(2a))$. Expanding and solving the resulting quadratic yields the result (negative root for stability).
\end{proof}

\begin{remark}[Light Separation]
$\kappa^2=\frac{1}{2}(3 - \sqrt{1 + 8\tau_W^2(O_o)})$ separates the gravitational part from the light signal.
\end{remark}

\subsubsection{Method~C: Arbitrary Phase Method}

\begin{theorem}[Arbitrary Phase Formula]
\hypertarget{thm:arbitrary_phase}{}\label{thm:arbitrary_phase}
\hypertarget{eq:Rs_arbitrary}{}\label{eq:Rs_arbitrary}
\begin{equation}
\boxed{R_s = \frac{r_o}{2(2a-r_o)}(4a - r_o - \sqrt{(4a-r_o)^2 - 8a(2a-r_o)(1-\tau_{Wo}^2)}).}
\end{equation}
\end{theorem}

\begin{proof}
Starting from $W = \frac{1}{2}(\kappa^2-\beta^2) = R_s/(4a)$, substituting into $\tau_{Wo}^2 = (1-\kappa^2)(1-\beta^2)$, and solving the resulting quadratic in $R_s$ yields the result.
\end{proof}

\subsection{The Role of $G$ as Translation Constant}

\begin{theorem}[Constants as Converters]
In WILL RG, $G$ and $m_0$ are derived calibration tools used to translate geometric scales into legacy units.
\end{theorem}

\begin{proof}
The operational procedure: (1)~Measure $\tau_W$ via spectroscopy; (2)~Measure $r$ via astrometry; (3)~Calculate $R_s = f(r,\tau_W)$. If one wishes to interface with legacy catalogues: $m_0 \equiv R_s c^2/(2G)$.
\end{proof}

\begin{remark}[Historical Artifact]
The kilogram is a human convention. The fundamental quantity is $\kappa^2 = R_s/r$, encoding the energy density ratio $\rho/\rho_{\max}$.
\end{remark}


% =============================================================================
% SECTION 21: ECCENTRICITY DERIVATION
% =============================================================================

\section{Derivation of Relational Eccentricity}
\hypertarget{sec:rel_ecc}{}\label{sec:rel_ecc}

\begin{theorem}[Geometric Eccentricity]
\begin{equation}
e = \frac{2\beta_p^2}{\kappa_p^2}-1=\frac{1}{\delta_p} - 1.
\end{equation}
\end{theorem}

\begin{proof}
From conservation of $W$ and the angular invariant $h = r\beta$ at the turning points. The apoapsis projections in terms of periapsis values: $\beta_a^2 = \beta_p^2((1-e)/(1+e))^2$ and $\kappa_a^2 = \kappa_p^2((1-e)/(1+e))$. Energy balance $W_p = W_a$ yields:
\[
\kappa_p^2[1 - \frac{1-e}{1+e}] = \beta_p^2[1 - (\frac{1-e}{1+e})^2],
\]
which reduces to $2\beta_p^2 = \kappa_p^2(1+e)$.
\hypertarget{eq:proj-balance}{}\label{eq:proj-balance}
Therefore $\delta_p = 1/(1+e)$ and
\[
\boxed{e = \frac{1}{\delta_p} - 1 = \frac{2\beta_p^2}{\kappa_p^2}-1}
\]
\end{proof}


% =============================================================================
% SECTION 22: PRECESSION
% =============================================================================

\section{Orbital Precession as State Difference Accumulation}
\hypertarget{sec:precession_law}{}\label{sec:precession_law}

The general precession law from geometric state difference accumulation:
\hypertarget{eq:precession_law}{}\label{eq:precession_law}
\begin{equation}
\Delta\varphi = \frac{2\pi\,Q^{2}}{1-e^{2}} = \frac{3\pi R_s}{a(1-e^2)}.
\end{equation}

Transforming to periapsis observables via $R_s = \kappa_p^2 r_p$ and $a(1-e^2) = r_p(1+e) = 2\beta_p^2 r_p/\kappa_p^2$:
\hypertarget{eq:light-ratio-law}{}\label{eq:light-ratio-law}
\begin{equation}
\boxed{\Delta\varphi = \frac{3}{2}\pi \frac{\kappa_p^4}{\beta_p^2}}
\end{equation}

\textbf{No differential equations. No metric. Pure algebra of light ratio.}

\subsection{Verification: Mercury}

$\kappa_p^4 \approx 4.11 \times 10^{-15}$, $\beta_p^2 \approx 3.87 \times 10^{-8}$:
\[
\Delta\varphi \approx \frac{3\pi}{2}(1.062 \times 10^{-7}) \approx 5.00 \times 10^{-7}\text{ rad/orbit} = \mathbf{43''}\text{/century}.
\]

\subsection{Verification: Star S2}

Data (GRAVITY Collaboration): $e \simeq 0.8846$, $\beta_p \simeq 0.0255$. Reconstructed $\kappa_p \approx 0.02627$. Prediction: $\Delta\varphi \approx 11.89'$. Observed: $12' \pm 1.5'$.


\subsection{Blind Prediction: S4716}

\begin{tcolorbox}[colback=blue!5, colframe=black!80!black, title=Blind Prediction Protocol]
Prediction made November 2025, prior to observational confirmation. De-projecting raw LOS velocity from 2009 SINFONI data yields $\Delta\varphi \approx \mathbf{14.80}$ arcmin/orbit.
\end{tcolorbox}

From $v_{\text{LOS}} \approx 1690$ km/s, de-projection factor $\mathcal{P} \approx 0.25$, yielding $\beta_o \approx 0.02255$. Propagating to periapsis via invariant $W$: $\beta_p \approx 0.0265$, $\kappa_p \approx 0.0283$.
\begin{equation}
\boxed{\Delta\varphi \approx 14.80\text{ arcmin/orbit}.}
\end{equation}

\begin{tcolorbox}[colback=gray!5, colframe=black!80!black, title=Desmos Project]
\textbf{\href{https://www.desmos.com/calculator/n4lmkpsebx}{DESMOS: R.O.M.}}
\end{tcolorbox}


% =============================================================================
% SECTION 23: CHRONO-SPECTROSCOPIC THEOREM
% =============================================================================

\section{Chrono-Spectroscopic Theorem: Absolute System Scale}
\hypertarget{sec:absolute_scale}{}\label{sec:absolute_scale}

For a complete orbital cycle, the absolute system scale decouples from spatial coordinates:
\[
R_{s} = T \cdot c \cdot \frac{\kappa^{2}\beta}{2\pi},
\]
generated exclusively from Chronometry (period $T$) and Spectroscopy (light shifts).

Validation on the Solar System via Mercury yields $R_s \approx 2953.3$~m, matching the classical $2GM_\odot/c^2$ without reference to $G$, $M$, or parallax.


% =============================================================================
% SECTION 24: M sin i DEGENERACY
% =============================================================================

\section{Sidestepping the Classical $M\sin i$ Degeneracy}
\hypertarget{sec:M_sin(i)}{}\label{sec:M_sin(i)}

R.O.M.\ sidesteps this degeneracy by introducing two independent higher-order kinematic invariants depending on $\beta_{\text{int}}$ alone:

\begin{theorem}[Decoupling of $\beta$ and $i$]
The classical linear amplitude $K \propto \beta_{\text{int}} \sin i$ is supplemented by:
\begin{enumerate}[nosep]
    \item \textbf{Precessional Phase Shift:} $\Delta\phi = 6\pi\beta_{\text{int}}^2$,
    \item \textbf{Systemic Transverse Baseline} $Z_{\text{sys}}(\phi)$.
\end{enumerate}
Both depend only on $\beta_{\text{int}}^2$, independent of $i$.
\end{theorem}

Once $\beta_{\text{int}}$ is fixed, the inclination is recovered from $\sin i = K\sqrt{1-e^2}/\beta_{\text{int}}$.

\subsection{Empirical Validation: MCMC on S0-2}
\hypertarget{sec:mcmc_s02}{}\label{sec:mcmc_s02}

Full R.O.M.\ successfully encompasses the GRAVITY Collaboration parameters within its 95\% CI, while the ablated (classical) model fails.

\begin{table}[H]
\centering\footnotesize
\renewcommand{\arraystretch}{1.1}
\begin{tabular}{|l|c|c|c|c|}
\hline
\textbf{Param.} & \textbf{Full R.O.M.} & \textbf{Ablated} & \textbf{GRAVITY} & \textbf{95\%?} \\
\hline
$\beta$ & $0.00597^{+0.00108}_{-0.00078}$ & $0.00477$ & $0.00645$ & \textbf{YES} \\
\hline
$i$ & $130.8^{+9.6}_{-11.7}$ & $105.8$ & $134.0$ & \textbf{YES} \\
\hline
$e$ & $0.88501$ & $0.88496$ & $0.88466$ & \textbf{YES} \\
\hline
\end{tabular}
\caption{MCMC posteriors for S0-2.}
\label{tab:mcmc_results}
\end{table}

\subsection{The Decryption Invariant}
\hypertarget{sec:analytical_sini}{}\label{sec:analytical_sini}

The inclination-free constraint:
\begin{equation}
Z_{rawmax} D_{max}(1 - e\cos\omega_i) + Z_{rawmin} D_{min}(1 + e\cos\omega_i) = 2,
\end{equation}
proves analytically that $\beta$ is determined by the asymmetry of transverse baselines, independent of viewing angle. Extraction:
\begin{equation}
\boxed{\sin i = \frac{\sqrt{1-e^2}}{2\beta}\left[ Z_{rawmax}D_{max} - Z_{rawmin}D_{min} \right]}
\end{equation}

\subsection{Synthetic Validation}
\hypertarget{sec:synthetic_validation}{}\label{sec:synthetic_validation}

Blind test on synthetic 1PN data: R.O.M.\ recovered true inclination $i = 49.92^\circ$ (true: $49.85^\circ$) from 1D radial velocity data alone.


% =============================================================================
% SECTION 25: GRAVITATIONAL OPTICS
% =============================================================================

\section{Gravitational Deflection and Lensing}
\hypertarget{sec:grav_optics}{}\label{sec:grav_optics}

\subsection{Unified Interaction Gradient}
\hypertarget{sec:grav_deflection}{}\label{sec:grav_deflection}

The geometric scaling factor $\Gamma$ is determined by the available phase buffer $\beta_Y$:
\begin{equation}
\Gamma(\beta) = \frac{1 + \beta^2}{2} = 1 - \frac{\beta_Y^2}{2}.
\end{equation}

The unified deflection for any trajectory:
\begin{equation}
\boxed{\Delta\varphi = 2\arcsin\!\left(\frac{\kappa_p^2(1+\beta_p^2)}{2\beta_p^2 - \kappa_p^2(1+\beta_p^2)}\right)}
\end{equation}

For light ($\beta_p=1$):
\begin{equation}
\boxed{\Delta\gamma = 2\arcsin\!\left(\frac{\kappa_p^2}{\kappa_{Xp}^2}\right)}
\end{equation}

\begin{tcolorbox}[colback=gray!5, colframe=black!80!black, title=Desmos Project]
\href{https://www.desmos.com/calculator/ldynwowqvi}{Algebraic Light Deflection}
\end{tcolorbox}

\subsection{Einstein Ring}
\hypertarget{sec:grav_lens}{}\label{sec:grav_lens}

\begin{theorem}[Exact Algebraic Einstein Ring]
For perfect alignment ($\theta_S = 0$):
\begin{equation}
\theta_E = 2\frac{D_{LS}}{D_S}\arcsin\!\left(\frac{\kappa_p^2(\theta_E)(1+\beta_p^2)}{2\beta_p^2 - \kappa_p^2(\theta_E)(1+\beta_p^2)}\right)
\end{equation}
\end{theorem}

In the photonic limit ($\beta_p\to 1$): $\theta_{E,\gamma} = \sqrt{2R_s D_{LS}/(D_L D_S)}$, matching full GR.


% =============================================================================
% SECTION 26: WILL INVARIANT
% =============================================================================

\section{$W_{\text{ILL}}$: Unity of Relational Structure}
\hypertarget{sec:willinvariant}{}\label{sec:willinvariant}

\begin{tcolorbox}[colback=gray!5, colframe=black!80!black, title=Desmos project]
\textbf{\href{https://www.desmos.com/calculator/wswnhckrru}{WILL = 1}}
\end{tcolorbox}

For any energy-closed system, WILL appears through four operational projections. In dimensionful form:
\begin{small}
\[
M = \frac{\beta^{2}}{\beta_{Y}}\frac{c^{2}a}{G}, \quad
E = \frac{\kappa^{2}}{\kappa_{X}}\frac{c^{4}a}{2G}, \quad
T = \kappa_{X}\!\left(\frac{2Gm_{0}}{\kappa^{2}c^{3}}\right)^{\!2}\!, \quad
L = \beta_{Y}\!\left(\frac{Gm_{0}}{\beta^{2}c^{2}}\right)^{\!2}\!.
\]
\end{small}

Combining:
\[
\boxed{W_{\text{ILL}} \equiv \frac{ET}{ML} = \frac{E_0 t^2}{m_0 a^2} = 1.}
\]

This holds for all energies, all scales, and all phases:
\[
\frac{E_{o}}{M_{o}} = \frac{L_{o}}{T_{o}},
\]
so the energy sector and spacetime sector are locked by a single relational constraint.

\[
\boxed{W_{\text{ILL}}= 1.}
\]


% =============================================================================
% SECTION 27: GENERATIVE PHYSICS
% =============================================================================

\section{Ontological Shift: From Descriptive to Generative Physics}
\hypertarget{Generative Physics}{}\label{Generative Physics}

In RG, laws are not added on top of observations; they are \emph{generated} as inevitable consequences of relational geometry. What classical physics calls ``laws of nature'' are secondary shadows of $\text{SPACETIME} \equiv \text{ENERGY}$.

\begin{table}[H]
\centering\footnotesize
\renewcommand{\arraystretch}{1.2}
\begin{tabular}{|p{3.5cm}|p{3.5cm}|}
\hline
\textbf{Descriptive (Standard)} & \textbf{Generative (WILL)} \\
\hline
Laws are assumptions that model reality &
Laws are identities enforced by geometry \\
\hline
Time/space are external backgrounds &
Time/space are projections of energy \\
\hline
Goal: describe what is observed &
Goal: show why nothing else is possible \\
\hline
\end{tabular}
\caption{Ontological contrast.}
\end{table}

\hypertarget{GRvsRG}{}\label{GRvsRG}
\begin{table}[H]
\centering\footnotesize
\renewcommand{\arraystretch}{1.1}
\begin{tabular}{|p{2cm}|p{2.5cm}|p{2.5cm}|}
\hline
\textbf{Phenomenon} & \textbf{GR} & \textbf{RG} \\
\hline
GPS time shift & Combination of SR+GR & $\tau=\beta_Y\kappa_X$ \\
\hline
Precession & Geodesic eqs.\ in Schwarzschild & $\Delta\varphi=\frac{3\pi}{2}\frac{\kappa_p^4}{\beta_p^2}$ \\
\hline
Singularities & $\rho\to\infty$ & Bounded by $\rho_{\max}$ \\
\hline
Local grav.\ energy & ``Cannot be localized'' & $\kappa$ directly measurable \\
\hline
\end{tabular}
\caption{GR vs.\ RG comparison.}
\end{table}


% =============================================================================
% SECTION 28: CONCLUSION
% =============================================================================

\section{Conclusion}

WILL Relational Geometry fully reproduces the predictive content of both Special and General Relativity while addressing their foundational inconsistencies: the lack of an operational definition of local gravitational energy density, the artificial separation of kinetic and gravitational energy, and the emergence of singularities as pathological artifacts.

From a single Ontological Principle---$\text{SPACETIME} \equiv \text{ENERGY}$---we derived the Lorentz factor, energy-momentum relation, gravitational time dilation, the closure law $\kappa^2=2\beta^2$, the Energy-Symmetry Law, orbital precession, gravitational deflection, the Einstein field equation, and the WILL invariant $W_{\text{ILL}}=1$.

\[
\boxed{\text{Special and General Relativity emerge from the same geometric principles.}}
\]

\begin{tcolorbox}[colback=gray!5, colframe=black!80!black, title=Final Summary]
\[
\textbf{SPACETIME} \;\equiv\; \textbf{ENERGY}.
\]
\end{tcolorbox}


% =============================================================================
% REFERENCES
% =============================================================================

\begin{thebibliography}{99}

\bibitem{Schwarzschild1916} Schwarzschild, K. (1916). \textit{Sitzungsberichte der K\"oniglich Preussischen Akademie der Wissenschaften}.

\bibitem{Tolman1939} Tolman, R.C. (1939). \textit{Phys.\ Rev.}\ \textbf{55}, 364; Oppenheimer, J.R., Volkoff, G.M. (1939). \textit{Phys.\ Rev.}\ \textbf{55}, 374.

\bibitem{Dyson1920} Dyson, F.W., Eddington, A.S., Davidson, C. (1920). \textit{Phil.\ Trans.\ R.\ Soc.\ A}\ \textbf{220}, 291--333.

\bibitem{Lebach1995} Lebach, D.E. et al. (1995). \textit{Phys.\ Rev.\ Lett.}\ \textbf{75}, 1439.

\bibitem{Ashby2003} Ashby, N. (2003). \textit{Living Rev.\ Relativity}\ \textbf{6}:1.

\bibitem{Will2014} Will, C.M. (2014). \textit{Living Rev.\ Relativity}\ \textbf{17}:4.

\bibitem{LLR} Williams, J.G., Turyshev, S.G., Boggs, D.H. (2006). \textit{Adv.\ Space Res.}\ \textbf{37}, 67--71.

\bibitem{Kramer2021} Kramer, M. et al. (2021). \textit{Phys.\ Rev.\ X}\ \textbf{11}, 041050.

\bibitem{Misner1973} Misner, Thorne, Wheeler (1973). \textit{Gravitation}.

\bibitem{Peissker2022} Pei{\ss}ker, F. et al. (2022). \textit{ApJ}\ \textbf{933}, 49.

\bibitem{NASA_Eclipse_Mercury} Espenak, F. (2014). Seven Century Catalog of Mercury Transits. \textit{NASA Eclipse Web Site}.

\bibitem{UniverseToday_Mercury} Williams, M. (2016). How Far is Mercury From the Sun? \textit{Universe Today}.

\bibitem{IAU_2015_ResB3} Mamajek, E.E. et al. (2015). IAU 2015 Resolution B3. \textit{arXiv:1510.07674}.

\end{thebibliography}

\end{multicols}

\end{document}
